% GENERATED FILE. Edit canonical lessons, then run npm run generate:books.
% canonical-source-hash: fnv1a64:c6612b92d14385ff
% canonical-lessons: TA-C453-attiram, TA-C453-ericcal, TA-C453-ciru, TA-C453-kontali, TA-C453-munumunu

\chapter{Rage and Irritation}
\label{ch:ta-rage-and-irritation}
\hlchaptermodality{453}

\begin{chapteropening}
\textbf{By the end of this chapter:} \emph{I can talk about rage, irritation and grumbling.}

Complete the last lesson of chapter 453: I can talk about rage, irritation and grumbling.
\end{chapteropening}

\section[\texorpdfstring{\ta{ஆத்திரம்} (āttiram)}{ஆத்திரம் (āttiram)}]{\ta{ஆத்திரம்} (āttiram) — rage; fury}

\label{lesson:TA-C453-attiram}



\begin{quote}
Before the new one: say the Tamil for a complaint, then the Tamil for an athlete.
\end{quote}

\subsection*{You'll want to know: \ta{ஆத்திரம்}}
\textbf{\ta{ஆத்திரம்}} — \emph{āttiram} — \textquotedblleft{}rage; fury\textquotedblright{}.

\begin{quote}
\ta{அவன்} \ta{ஆத்திரத்தில்} \ta{கத்தினான்}.
\end{quote}

\begin{quote}
\emph{avaṉ āttirattil kattiṉāṉ.}
\end{quote}

\begin{quote}
He shouted in a rage.
\end{quote}

A -\ta{அம்} noun takes the oblique -\ta{அத்து}- before an ending. In a rage: \ta{ஆத்திரத்தில்}.

\begin{cousinweb}
The origin of \textbf{\ta{ஆத்திரம்}} (\emph{āttiram}) is not settled. It is often set beside Sanskrit \emph{ātura}, \textquotedblleft{}suffering, distressed\textquotedblright{} (Monier-Williams), but the link is uncertain.
\end{cousinweb}

\begin{grammarlens}[title={What you've built}]
One more word for describing what you lived and felt.
\end{grammarlens}

\subsection*{Your turn}
\begin{itemize}
\raggedright
  \item \emph{Say it:} \emph{āttiram}
  \item \emph{Say it:} \emph{āttiram}, once more
  \item \emph{Say it:} \emph{avaṉ āttirattil kattiṉāṉ.} — He shouted in a rage.
  \item \emph{Recall:} say the Tamil for a complaint, then the Tamil for an athlete, then say \emph{āttiram} again
\end{itemize}

\begin{tcolorbox}[breakable,colback=teal!4,colframe=teal!35!black,title={Before you move on}]
What does \ta{ஆத்திரம்} mean? (\textquotedblleft{}Rage; fury\textquotedblright{}.) Say it once more.
\end{tcolorbox}
\section[\texorpdfstring{\ta{எரிச்சல்} (ericcal)}{எரிச்சல் (ericcal)}]{\ta{எரிச்சல்} (ericcal) — irritation; a burning feeling}

\label{lesson:TA-C453-ericcal}



\begin{quote}
Before the new one: say the Tamil for an athlete, then the Tamil for rage.
\end{quote}

\subsection*{You'll want to know: \ta{எரிச்சல்}}
\textbf{\ta{எரிச்சல்}} — \emph{ericcal} — \textquotedblleft{}irritation; a burning feeling\textquotedblright{}.

\begin{quote}
\ta{சத்தம்} \ta{கேட்டு} \ta{எனக்கு} \ta{எரிச்சல்} \ta{வந்தது}.
\end{quote}

\begin{quote}
\emph{cattam kēṭṭu eṉakku ericcal vantatu.}
\end{quote}

\begin{quote}
The noise irritated me.
\end{quote}

A feeling comes to the person in the dative: \ta{எனக்கு} \ta{எரிச்சல்} \ta{வந்தது}, literally irritation came to me. Of the body too: \ta{கண்} \ta{எரிச்சல்}, stinging eyes.

\begin{cousinweb}
\textbf{\ta{எரிச்சல்}} (\emph{ericcal}) is native Dravidian, a noun from \ta{எரி} \emph{eri} \textquotedblleft{}to burn\textquotedblright{} (DEDR 811). The same entry gives the verb the sense of smarting pain and painful feeling.
\end{cousinweb}

\begin{grammarlens}[title={What you've built}]
One more word for describing what you lived and felt.
\end{grammarlens}

\subsection*{Your turn}
\begin{itemize}
\raggedright
  \item \emph{Say it:} \emph{ericcal}
  \item \emph{Say it:} \emph{ericcal}, once more
  \item \emph{Say it:} \emph{cattam kēṭṭu eṉakku ericcal vantatu.} — The noise irritated me.
  \item \emph{Recall:} say the Tamil for an athlete, then the Tamil for rage, then say \emph{ericcal} again
\end{itemize}

\begin{tcolorbox}[breakable,colback=teal!4,colframe=teal!35!black,title={Before you move on}]
What does \ta{எரிச்சல்} mean? (\textquotedblleft{}Irritation; a burning feeling\textquotedblright{}.) Say it once more.
\end{tcolorbox}
\section[\texorpdfstring{\ta{சீறு} (cīṟu)}{சீறு (cīṟu)}]{\ta{சீறு} (cīṟu) — to flare up; to hiss}

\label{lesson:TA-C453-ciru}



\begin{quote}
Before the new one: say the Tamil for rage, then the Tamil for irritation.
\end{quote}

\subsection*{You'll want to know: \ta{சீறு}}
\textbf{\ta{சீறு}} — \emph{cīṟu} — \textquotedblleft{}to flare up; to hiss\textquotedblright{}.

\begin{quote}
\ta{நான்} \ta{கேள்வி} \ta{கேட்டேன்}, \ta{அவர்} \ta{சீறினார்}.
\end{quote}

\begin{quote}
\emph{nāṉ kēḷvi kēṭṭēṉ, avar cīṟiṉār.}
\end{quote}

\begin{quote}
I asked a question and he flared up.
\end{quote}

Past (I …): \ta{சீறினேன்}, \emph{cīṟiṉēṉ}.

Its first sense is the hiss of a snake or the snort of a horse; of a person, a sudden angry outburst.

\begin{cousinweb}
\textbf{\ta{சீறு}} (\emph{cīṟu}) is native Dravidian (DEDR 2639): the entry joins the hiss of a snake, the snort of a horse and a burst of rage.
\end{cousinweb}

\begin{grammarlens}[title={What you've built}]
One more word for describing what you lived and felt.
\end{grammarlens}

\subsection*{Your turn}
\begin{itemize}
\raggedright
  \item \emph{Say it:} \emph{cīṟu}
  \item \emph{Say it:} \emph{cīṟu}, once more
  \item \emph{Say it:} \emph{nāṉ kēḷvi kēṭṭēṉ, avar cīṟiṉār.} — I asked a question and he flared up.
  \item \emph{Recall:} say the Tamil for rage, then the Tamil for irritation, then say \emph{cīṟu} again
\end{itemize}

\begin{tcolorbox}[breakable,colback=teal!4,colframe=teal!35!black,title={Before you move on}]
What does \ta{சீறு} mean? (\textquotedblleft{}To flare up; to hiss\textquotedblright{}.) Say it once more.
\end{tcolorbox}
\section[\texorpdfstring{\ta{கொந்தளி} (kontaḷi)}{கொந்தளி (kontaḷi)}]{\ta{கொந்தளி} (kontaḷi) — to seethe; to be in turmoil}

\label{lesson:TA-C453-kontali}



\begin{quote}
Before the new one: say the Tamil for irritation, then the Tamil for to flare up.
\end{quote}

\subsection*{You'll want to know: \ta{கொந்தளி}}
\textbf{\ta{கொந்தளி}} — \emph{kontaḷi} — \textquotedblleft{}to seethe; to be in turmoil\textquotedblright{}.

\begin{quote}
\ta{அந்தச்} \ta{செய்தியைக்} \ta{கேட்டு} \ta{அவன்} \ta{மனம்} \ta{கொந்தளித்தது}.
\end{quote}

\begin{quote}
\emph{antac ceytiyaik kēṭṭu avaṉ maṉam kontaḷittatu.}
\end{quote}

\begin{quote}
Hearing that news, he was in turmoil.
\end{quote}

Past (I …): \ta{கொந்தளித்தேன்}, \emph{kontaḷittēṉ}.

Also of a rough sea: \ta{கடல்} \ta{கொந்தளித்தது}, the sea ran high.

\begin{cousinweb}
\textbf{\ta{கொந்தளி}} (\emph{kontaḷi}) is native Dravidian (DEDR 2098); the same entry has \ta{கொந்தளம்} \emph{kontaḷam}, \textquotedblleft{}turmoil, confusion\textquotedblright{}.
\end{cousinweb}

\begin{grammarlens}[title={What you've built}]
One more word for describing what you lived and felt.
\end{grammarlens}

\subsection*{Your turn}
\begin{itemize}
\raggedright
  \item \emph{Say it:} \emph{kontaḷi}
  \item \emph{Say it:} \emph{kontaḷi}, once more
  \item \emph{Say it:} \emph{antac ceytiyaik kēṭṭu avaṉ maṉam kontaḷittatu.} — Hearing that news, he was in turmoil.
  \item \emph{Recall:} say the Tamil for irritation, then the Tamil for to flare up, then say \emph{kontaḷi} again
\end{itemize}

\begin{tcolorbox}[breakable,colback=teal!4,colframe=teal!35!black,title={Before you move on}]
What does \ta{கொந்தளி} mean? (\textquotedblleft{}To seethe; to be in turmoil\textquotedblright{}.) Say it once more.
\end{tcolorbox}
\section[\texorpdfstring{\ta{முணுமுணு} (muṇumuṇu)}{முணுமுணு (muṇumuṇu)}]{\ta{முணுமுணு} (muṇumuṇu) — to mutter; to grumble}

\label{lesson:TA-C453-munumunu}



\begin{quote}
Before the new one: say the Tamil for to flare up, then the Tamil for to seethe.
\end{quote}

\subsection*{You'll want to know: \ta{முணுமுணு}}
\textbf{\ta{முணுமுணு}} — \emph{muṇumuṇu} — \textquotedblleft{}to mutter; to grumble\textquotedblright{}.

\begin{quote}
\ta{அவன்} \ta{வேலை} \ta{செய்யும்போது} \ta{முணுமுணுத்தான்}.
\end{quote}

\begin{quote}
\emph{avaṉ vēlai ceyyumpōtu muṇumuṇuttāṉ.}
\end{quote}

\begin{quote}
He grumbled while he worked.
\end{quote}

Past (I …): \ta{முணுமுணுத்தேன்}, \emph{muṇumuṇuttēṉ}.

A doubled sound word, like many Tamil verbs of noise. Past with -\ta{த்த}-: \ta{முணுமுணுத்தான்}.

\begin{cousinweb}
\textbf{\ta{முணுமுணு}} (\emph{muṇumuṇu}) is native Dravidian (DEDR 4856), a doubled word that imitates a low, repeated sound.
\end{cousinweb}

\begin{grammarlens}[title={What you've built}]
One more word for describing what you lived and felt.
\end{grammarlens}

\subsection*{Your turn}
\begin{itemize}
\raggedright
  \item \emph{Say it:} \emph{muṇumuṇu}
  \item \emph{Say it:} \emph{muṇumuṇu}, once more
  \item \emph{Say it:} \emph{avaṉ vēlai ceyyumpōtu muṇumuṇuttāṉ.} — He grumbled while he worked.
  \item \emph{Recall:} say the Tamil for to flare up, then the Tamil for to seethe, then say \emph{muṇumuṇu} again
\end{itemize}

\begin{tcolorbox}[breakable,colback=teal!4,colframe=teal!35!black,title={Before you move on}]
What does \ta{முணுமுணு} mean? (\textquotedblleft{}To mutter; to grumble\textquotedblright{}.) Say it once more.
\end{tcolorbox}
